\clearpage
\SourceLeaf{457}
\begin{center}\large SIMPLEST DERIVATION OF THE\\FUNDAMENTAL THEOREM\end{center}
\begin{center}CONCERNING\end{center}
\begin{center}\large THE SHORTEST LINES ON SURFACES OF REVOLUTION.\end{center}
\begin{center}\begin{tabular}{lp{.73\linewidth}}
$P,P'$ & two points on such a surface\\
$r$ & length of the shortest line between $P$ and $P'$\\
$A,A'$ & its azimuths at $P$ and $P'$,\\
& the latter reckoned from east to west\\
$\varphi,\lambda;\ \varphi',\lambda'$ & latitude and longitude of the points $P,P'$\\
$\rho,\rho'$ & radii of the parallels through $P$ and $P'$\\
$R,R'$ & radii of curvature of the meridians at these points.
\end{tabular}\end{center}
\noindent $r+dr,\ \varphi+d\varphi,\ \lambda+d\lambda,\ \varphi'+d\varphi',\ \lambda'+d\lambda'$ are the values taken by $r,\varphi,\lambda,\varphi',\lambda'$ when two other points infinitely close to them are substituted for $P,P'$.

We then have
\[
dr=R\cos A.d\varphi-R'\cos A'.d\varphi'-\rho\sin A.d\lambda+\rho'\sin A'.d\lambda'.
\]
\noindent But evidently $dr=0$ when the conditions
\[
d\varphi=0,\qquad d\varphi'=0,\qquad d\lambda'=d\lambda
\]
\noindent are imposed. Thus
\[
\rho\sin A=\rho'\sin A'.
\]
\begin{flushright}Q. E. D.\end{flushright}
\Needspace{7\baselineskip}
\begin{center}\textbf{REMARK.}\end{center}
\noindent\hspace*{1.2em}The preceding note, which is found in a handbook, probably dates from the early forties.
\begin{flushright}\textsc{Stäckel}.\end{flushright}

